\subsubsection{Buffer Caches}
To be quick (as covered in SPCA in Semester 3), you need data to be available in memory, and even better still in the CPU caches.

This is no different for database systems.
Here we are however not dealing with the hardware side of caching, but caching commonly used pages and metadata in main memory,
as opposed to streaming them from disk on each access. However, database systems often write dirty blocks back to disk when dirty (i.e. use write-through caches),
or when in need of space.

Thus, in principle, buffer caches are similar to OS virtual memory and paging mechanisms, however, unlike OSes,
the database system knows the access patterns and can (thus) optimize much more.
